\section{Introducción: por qué la capacidad de razonamiento es clave para la IA}

Denny Zhou es científico investigador en Google DeepMind, dedicado desde hace mucho tiempo a la investigación de la capacidad de razonamiento de los modelos de lenguaje grandes. Es uno de los contribuyentes centrales de trabajos histĂłricos como Chain-of-Thought Prompting y Self-Consistency. Esta conferencia parte de una pregunta fundamental —\textbf{¿qué es realmente la inteligencia?}— para desarrollar el hilo central del razonamiento en los LLM.

\begin{knowledgebox}{El cambio de paradigma del aprendizaje automático al razonamiento}
El aprendizaje automático tradicional (few-shot learning, active learning, meta-learning, etc.) persigue la eficiencia de datos (data efficiency), pero en la práctica estos métodos no logran realmente «aprender a partir de pocas muestras». Denny Zhou sostiene que la razón fundamental por la que los seres humanos pueden aprender a partir de pocas muestras no son las regularidades estadísticas, sino la \textbf{capacidad de razonamiento (reasoning)}. Esta intuición impulsó su transición del aprendizaje automático hacia la investigación del razonamiento en los LLM.
\end{knowledgebox}

\subsection{Last Letter Concatenation: un ejemplo simple pero profundo}

Denny Zhou utiliza la tarea de Last Letter Concatenation para revelar la enorme diferencia entre los métodos tradicionales y los métodos de razonamiento:

\begin{itemize}
    \item \textbf{Definición de la tarea}: dado el nombre de una persona (por ejemplo, ``Elon Musk''), producir la concatenación de la última letra del nombre y la última letra del apellido (``n'' + ``k'' = ``nk'').
    \item \textbf{Método tradicional de aprendizaje automático}: entrenado con un Transformer, requiere miles de muestras etiquetadas y alcanza una exactitud de aproximadamente 85--90\%.
    \item \textbf{Few-shot Prompting}: al darle directamente al LLM unos cuantos ejemplos, este sigue cometiendo errores con facilidad (por ejemplo, para ``Barack Obama'' produce ``ck'' en lugar de ``ka'').
    \item \textbf{Chain-of-Thought Prompting}: al agregar pasos de razonamiento dentro de los ejemplos (``The last letter of Elon is n, the last letter of Musk is k, so nk''), basta con un solo ejemplo para alcanzar una exactitud de 100\%.
\end{itemize}

\begin{importantbox}{Intuición central}
Para una tarea que resulta simple para los seres humanos, si un método necesita una cantidad masiva de datos etiquetados para aprenderla, entonces es difícil considerarlo verdadera «inteligencia». La capacidad de razonamiento permite que los LLM generalicen a partir de muy pocas muestras, algo que los métodos tradicionales de aprendizaje automático no pueden igualar.
\end{importantbox}

\subsection{Resumen de la sección}

La capacidad de razonamiento es la clave para cerrar la brecha entre el «aprendizaje impulsado por datos» y la «inteligencia real». Al generar pasos intermedios al estilo Chain-of-Thought, los LLM logran un salto cualitativo: pasar de necesitar grandes cantidades de datos de entrenamiento a requerir solo unas cuantas muestras, o incluso ninguna.
